%\begin{name}
%	{Biên soạn: Thầy Lê Hồng Phi  \& Phản biện: Thầy Nguyễn Thành Nhân}
%	{Đề chuẩn hóa}
%\end{name}

\caulc
\Opensolutionfile{ans}[ans/ans-cau-TN-1]

\begin{ex}%[Dự án chuẩn hóa 10-11, cấu trúc Bộ 2025]%[Lê Hồng Phi]%[1D9N2-3]
	Do con súc sắc không cân đối nên mỗi khi gieo con súc sắc đó thì xác suất xuất hiện mặt $2$ chấm và mặt $5$ chấm lần lượt là $0{,}35$ và $0{,}25$. Xác suất xuất hiện mặt $2$ chấm hoặc mặt $5$ chấm là
	\choice
	{\True $0{,}6$}
	{$0{,}0875$}
	{$0{,}4875$}
	{$0{,9125}$}
	\loigiai{
		Gọi $A$ là biến cố xuất hiện mặt $2$ chấm khi gieo con súc sắc.\\
		Gọi $B$ là biến cố xuất hiện mặt $5$ chấm khi gieo con súc sắc.\\
		Do đó, $A\cup B$ là biến cố xuất hiện mặt $2$ chấm hoặc mặt $5$ chấm.\\
Ta có $\mathrm{P}(A)=0{,}35$ và $\mathrm{P}(B)=0{,}25$.\\
Do $A$ và $B$ là hai biến cố xung khắc nhau nên
\[\mathrm{P}(A\cup B)=\mathrm{P}(A)+\mathrm{P}(B)=0{,}35+0{,}25=0{,}6.\]
	}
\end{ex}

\begin{ex}%[Dự án chuẩn hóa 10-11, cấu trúc Bộ 2025]%[Lê Hồng Phi]%[1D5N1-2]
\immini[thm]{Thống kê chiều cao của học sinh lớp $11 \mathrm{A}$ ta có bảng số liệu như hình bên. Hỏi lớp có bao nhiêu học sinh có chiều cao từ $168 \mathrm{~cm}$ trở lên?
\choice
{\True$11$}
{$20$}
{$31$}
{$8$}}{\begin{tabular}{|c|c|}
\hline
Chiều cao (cm) & Số học sinh\\
\hline
$[150 ; 156)$ & $8$\\
$[156 ; 162)$ & $12$\\
$[162 ; 168)$ & $11$\\
$[168 ; 174)$ & $8$\\
$[174 ; 180)$ & $3$\\
\hline
\end{tabular}}
\loigiai{
Số học sinh có chiều cao từ $168 \mathrm{~cm}$ trở lên là $8+3=11$.
	}
\end{ex}


\begin{ex}%[Dự án chuẩn hóa 10-11, cấu trúc Bộ 2025]%[Lê Hồng Phi]%[1D9H2-3]
	Một hộp chứa $6$ quả cầu trắng và $7$ quả cầu đen. Lấy ngẫu nhiên đồng thời bốn quả trong số đó. Xác suất để có ít nhất một quả cầu trắng bằng
	\choice
	{$\dfrac{7}{143}$}
	{$\dfrac{120}{143}$}
	{$\dfrac{13}{24}$}
	{\True $\dfrac{136}{143}$}
	\loigiai{
		Số phần tử của không gian mẫu là $n(\Omega)=\mathrm{C}_{13}^4$.\\
		Gọi $A$ là biến cố lấy được ít nhất một quả cầu trắng.\\
		Biến cố đối của $A$ là biến cố $B$ \lq\lq không lấy được quả cầu trắng nào\rq\rq.\\
		Số phần tử của biến cố $B$ là $n(B)=\mathrm{C}_7^4$.\\
		Xác suất $\mathrm{P}(B)=\dfrac{n(B)}{n(\Omega)}=\dfrac{\mathrm{C}_7^4}{\mathrm{C}_{13}^4}=\dfrac{7}{143}$.\\
		Suy ra xác suất $\mathrm{P}(A)=1-\mathrm{P}(B)=\dfrac{136}{143}$.
	}
\end{ex}

\begin{ex}%[Dự án chuẩn hóa 10-11, cấu trúc Bộ 2025]%[Lê Hồng Phi]%[1D6N1-2]
	Cho $x, y$ là hai số thực dương và $m, n$ là hai số thực tùy ý. Đẳng thức nào sau đây là {\bfseries sai}?
	\choice
	{$\left(x^{m}\right)^{n}=(x)^{m n}$}
	{$x^{m}\cdot x^{n}=x^{m+n}$}
	{$(x y)^{n}=x^{n}y^{n}$}
	{\True$x^{m}y^{n}=(x y)^{m+n}$}
	\loigiai{Đẳng thức $x^{m}y^{n}=(x y)^{m+n}$ sai.
	}
\end{ex}

\begin{ex}%[Dự án chuẩn hóa 10-11, cấu trúc Bộ 2025]%[Lê Hồng Phi]%[1D6H2-2]
	Với $a$ và $b$ là các số thực dương và $a \neq 1$. Biểu thức $\log _{a}\left(a^2 b\right)$ bằng
	\choice
	{$2 \log _{a}b$}
	{$1+2 \log _{a}b$}
	{$2-\log _{a}b$}
	{\True $2+\log _{a}b$}
\loigiai{
	Ta có $\log _{a}\left(a^2 b\right)= \log _{a} a^2 + \log _{a} b=2\log _{a} a+\log _{a}b=2+\log _{a}b$.
}
\end{ex}

\begin{ex}%[Dự án chuẩn hóa 10-11, cấu trúc Bộ 2025]%[Lê Hồng Phi]%[1D7N2-1]
	Khẳng định nào sau đây là \textbf{sai}?
	\choice
	{$(\sin x)'=\cos x$}
	{$(\cot x)'=\dfrac{-1}{\sin^2x}$}
	{\True $(\tan x)'=\dfrac{1}{\sin^2x}$}
	{$(\cos x)'=-\sin x$}
	\loigiai{
		Theo công thức tính đạo hàm của hàm số lượng giác, khẳng định $(\tan x)'=\dfrac{1}{\sin^2x}$ là sai. Công thức đúng là $(\tan x)'=\dfrac{1}{\cos^2x}$.
	}
\end{ex}

\begin{ex}%[Dự án chuẩn hóa 10-11, cấu trúc Bộ 2025]%[Lê Hồng Phi]%[1D7H2-1]
	Tính đạo hàm của hàm số $y=\dfrac{2x}{x-1}$.
	\choice
	{$y'=\dfrac{-1}{(x-1)^2}$}
	{$y'=\dfrac{2}{(x-1)^2}$}
	{$y'=\dfrac{3}{(x-1)^2}$}
	{\True $y'=\dfrac{-2}{(x-1)^2}$}
	\loigiai{
		Ta có $y'=\dfrac{-2}{(x-1)^2}$.
	}
\end{ex}

\begin{ex}%[Dự án chuẩn hóa 10-11, cấu trúc Bộ 2025]%[Lê Hồng Phi]%[1D6H3-2]
	Tìm tập xác định của hàm số $y=\ln \left(-2x^2+x+3\right)$.
	\choice
	{$\mathscr{D}=(-\infty;-1) \cup\left(\dfrac{3}{2};+\infty\right)$}
	{$\mathscr{D}=\left[-1; \dfrac{3}{2}\right]$}
	{$\mathscr{D}=(-\infty;-1] \cup\left[\dfrac{3}{2};+\infty\right)$}
	{\True $\mathscr{D}=\left(-1; \dfrac{3}{2}\right)$}
	\loigiai{
		Hàm số xác định khi $-2x^2+x+3>0\Leftrightarrow -1<x<\dfrac{3}{2}$.\\
		Vậy tập xác định của hàm số là $\mathscr{D}=\left(-1; \dfrac{3}{2}\right)$.
	}
\end{ex}

\begin{ex}%[Dự án chuẩn hóa 10-11, cấu trúc Bộ 2025]%[Lê Hồng Phi]%[1D7H2-1]
Tính đạo hàm của hàm số $f(x)=\dfrac{1}{x}$ tại điểm $x=\dfrac{1}{2}$.
\choice
{$f'(\frac{1}{2})=-\dfrac{1}{4}$}
{$f'(\frac{1}{2})=\dfrac{1}{4}$}
{\True $f'(\frac{1}{2})=-4$}
{$f'(\frac{1}{2})=4$}
\loigiai{
Ta có $f'(x)=-\dfrac{1}{x^2}$, thay $x=\dfrac{1}{2}$ ta được $f'(\frac{1}{2})=-4$.
}
\end{ex}

\begin{ex}%[Dự án chuẩn hóa 10-11, cấu trúc Bộ 2025]%[Lê Hồng Phi]%[1H8H6-2]
	\immini
	{	Trong hình bên, máy tính xách tay đang mở gợi nên hình ảnh của một góc nhị diện. Ta gọi số đo góc nhị diện đó là độ mở của màn hình máy tính. Tính độ mở của màn hình máy tính theo đơn vị độ, biết tam giác $ABC$ có độ dài các cạnh là $AB=AC=30$ cm và $BC=30\sqrt{3}$ cm.
		\choice
		{$90^\circ$}
		{$60^\circ$}
		{$45^\circ$}
		{\True $120^\circ$}
	}
	{
		\begin{tikzpicture}[scale=0.7, font=\footnotesize, line join=round, line cap=round, >=stealth]
			\path
			(0,0) coordinate (A)
			(-1,5) coordinate (B)
			(3,4) coordinate (D)
			(4,1) coordinate (E)
			(3.5,-1) coordinate (C)
			(7,0.2) coordinate (F)
			($(A)+(0.5,0)$) coordinate (A')
			($(A)!10/11!(E)$) coordinate (M)
			($(M)+(0.5,0)$) coordinate (M')
			($(A)!1/2!(C)$) coordinate (N)
			($(N)+(0.5,0)$) coordinate (N')
			($(E)!0.65!(F)$) coordinate (O)
			($(O)-(0.3,0)$) coordinate (O')
			;
			\draw[blue!50,line width=2pt] (A)--(C)--(F);
			\draw[blue!70,line width=3pt] (A)--(B)--(D)--(E)--(A)
			;
			\fill[blue!30] (O')--(N')--(A')--(M')--cycle;
			\draw(F)--(E) (B)--(C);
			\foreach \x/\g in {A/160,B/180,C/270} \fill[black] (\x) circle (2pt)+(\g:0.6) node{$\x$};
		\end{tikzpicture}
	}
	\loigiai{
		Góc nhị diện đã cho có số đo bằng số đo của góc $\widehat{BAC}$.\\
		Trong tam giác $ABC$ ta có $\cos \widehat{BAC}=\dfrac{AB^2+AC^2-BC^2}{2\cdot AB \cdot AC}=\dfrac{30^2+30^2-\left(30\sqrt{3}\right)^2}{2\cdot 30 \cdot 30}=-\dfrac{1}{2}\Rightarrow \widehat{BAC}=120^\circ$.\\
		Vậy độ mở của màn hình máy tính là $120^\circ$.
	}
\end{ex}

\begin{ex}%[Dự án chuẩn hóa 10-11, cấu trúc Bộ 2025]%[Lê Hồng Phi]%[1H8H2-2]
	\immini{Cho hình chóp $S.ABCD$  có đáy $ABCD$ là hình chữ nhật và $SA$ vuông góc với mặt phẳng $(ABCD)$. Mặt phẳng nào dưới đây vuông góc với đường thẳng $BC$?
		\choice
		{$(SBD)$}
		{\True $(SAB)$}
		{$(SAC)$}
		{$(SCD)$}}{\begin{tikzpicture}[scale=1, font=\footnotesize, line join=round, line cap=round,>=stealth]
\path (0,0)coordinate (B)++(3,0) coordinate (C)++(1.3,1) coordinate (D) ($(B)+(D)-(C)$) coordinate (A) (0,2.5) coordinate (h) ($(A)+(h)$) coordinate (S) ($(A)!0.3!(C)$) coordinate (H);
\draw (S)--(B)--(C)--(D)--cycle (S)--(C);
\draw[dashed] (S)--(A)--(B) (C)--(A)--(D)--(B);
\foreach \p/\g in {A/150,B/240,C/-60, D/0, S/90} \fill[black] (\p) circle(1pt)+(\g:0.3) node{$\p$};
%\foreach \x/\y/\z in {S/A/B,S/A/D,A/B/C,B/C/D,C/D/A,D/A/B}\draw pic[draw, angle radius=0.2cm]{right angle=\x--\y--\z};
\end{tikzpicture}
}
	\loigiai{
		Vì $ABCD$ là hình chữ nhật nên $BC\perp AB$. (1)\\
		Mặt khác $SA\perp (ABCD)\Rightarrow BC\perp SA$. (2) \\
		Từ $(1)$ và $(2)$ suy ra $BC\perp (SAB).$	
	}
\end{ex}

\begin{ex}%[Dự án chuẩn hóa 10-11, cấu trúc Bộ 2025]%[Lê Hồng Phi]%[1H8H4-2]
	Cho hình chóp $S.ABCD$ có đáy $ABCD$ là hình chữ nhật, cạnh bên $SA$ vuông góc với đáy. Khẳng định nào sau đây \textbf{sai}?
	\choice
	{\True $(SBD)\perp (SAC)$}
	{$(SCD)\perp (SAD)$}
	{$(SAC)\perp (ABC)$}
	{$(SBC)\perp (SAB)$}
	\loigiai{
		\immini{
			Gọi $H$ là hình chiếu vuông góc của $B$ trên $AC$ (Do $ABCD$ là hình chữ nhật và không là hình vuông nên $H$ không trùng với tâm của hình hình chữ nhật $ABCD$).\\ 
			Khi đó, $BH\perp (SAC)\Rightarrow (SBH)\perp (SAC)$.\\
			Vậy nếu $(SBD)\perp (SAC)$ thì giao tuyến của $(SBD)$ và $(SBH)$ là $SB$ phải vuông góc với $(SAC)$. Điều này không xảy ra vì $S.ABCD$ là hình chóp.\\
			Vậy khẳng định $(SBD)\perp (SAC)$ là sai. 
		}
		{\begin{tikzpicture}[scale=1, font=\footnotesize, line join=round, line cap=round,>=stealth]
\path (0,0)coordinate (B)++(4,0) coordinate (C)++(1.5,1.3) coordinate (D) ($(B)+(D)-(C)$) coordinate (A) (0,2.5) coordinate (h) ($(A)+(h)$) coordinate (S) ($(A)!0.3!(C)$) coordinate (H);
\draw (S)--(B)--(C)--(D)--cycle (S)--(C);
\draw[dashed] (S)--(A)--(B)--(H) (C)--(A)--(D)--(B);
\foreach \p/\g in {A/150,B/240,C/-60, D/0, S/90,H/45} \fill[black] (\p) circle(1pt)+(\g:0.3) node{$\p$};
\foreach \x/\y/\z in {S/A/B,S/A/D,B/H/C,A/B/C,B/C/D,C/D/A,D/A/B}\draw pic[draw, angle radius=0.2cm]{right angle=\x--\y--\z};
\end{tikzpicture}
}
	}
\end{ex}
\Closesolutionfile{ans}
\indapan{7}{ans/ans-cau-TN-1}

\cauds
\Opensolutionfile{ans}[ans/ans-cau-DS-1]
\begin{ex}%[Dự án chuẩn hóa 10-11, cấu trúc Bộ 2025]%[Lê Hồng Phi]%[1D9H2-1]
	Một hộp có $20$ chiếc thẻ cùng loại, mỗi thẻ được ghi một số trong các số {$1,2,\ldots,19,20$}; hai thẻ khác nhau thì ghi hai số khác nhau. Rút ngẫu nhiên một chiếc thẻ trong hộp. Xét các biến cố\\
	\indent $A\colon$ \lq\lq Số trên thẻ được rút ra là số chia hết cho $5$\rq\rq.\\
	\indent $B\colon$ \lq\lq Số trên thẻ được rút ra là số chia hết cho $3$\rq\rq.\\
	\indent $C\colon$ \lq\lq Số trên thẻ được rút ra là số chia hết cho $5$ hoặc chia hết cho $3$\rq\rq.\\
	\indent $D\colon$ \lq\lq Số trên thẻ được rút ra là số chia hết cho $15$\rq\rq.
	\choiceTF
	{Biến cố $A$ và biến cố $D$ là hai biến cố xung khắc}
	{\True Xác suất của biến cố $D$ là $\dfrac{1}{20}$}
	{Biến cố $D$ là biến cố hợp của biến cố $A$ và biến cố $B$}
	{Số phần tử của không gian mẫu là $190$}
	\loigiai{
		\begin{itemchoice}
			\itemch Ta có thẻ số $15$ vừa chia hết cho cả $3$ và $5$ nên $15$ vừa là phần tử của biến cố $A$, vừa là phần tử của biến cố $B$ nên hai biến cố $A$, $B$ không xung khắc.
			\itemch Xác suất của biến cố $D$ là $\mathrm{P}(D)=\dfrac{1}{20}$.
			\itemch Ta có $D=A\cap B$.
			\itemch Số phần tử của không gian mẫu là $n(\Omega)=\mathrm{C}_{20}^1=20$.
		\end{itemchoice}
	}
\end{ex}

\begin{ex}%[Dự án chuẩn hóa 10-11, cấu trúc Bộ 2025]%[Lê Hồng Phi]%[1D6V4-5]
	Cho hàm số $f(x)=\log_3\left(3^{x+1}-1\right)$.
	\choiceTF
	{\True Biết nghiệm của phương trình $f(x)=x$ là $x=\log_ab$. Giá trị của $\dfrac{a}{b}$ là $6$}
	{Nghiệm của phương trình $f(x)=0$ là $x=1-\log_32$}
	{\True Điều kiện xác định của hàm số $f(x)$ là $x>-1$}
	{Tập nghiệm của bất phương trình $f(x)<0$ là $S=(-1;+\infty)$}
	\loigiai{
		\begin{itemchoice}
			\itemch Ta có
			$\log_3\left(3^{x+1}-1\right)=x\Leftrightarrow3^{x+1}-1=3^x=\Leftrightarrow2\cdot3^x=1\Leftrightarrow x=\log_3\dfrac{1}{2}.$
			Khi đó $a=3$ và $b=\dfrac{1}{2}$, nên $\dfrac{a}{b}=6$.
			\itemch Ta có
			$\log_3\left(3^{x+1}-1\right)=0\Leftrightarrow3^{x+1}-1=1\Leftrightarrow3^{x+1}=2\Leftrightarrow x=\log_32-1.$
			\itemch Điều kiện
			$3^{x+1}-1>0\Leftrightarrow3^{x+1}>1\Leftrightarrow x+1>0\Leftrightarrow x>-1.$
			\itemch Ta có
			$f(x)<0\Leftrightarrow3^{x+1}-1<0\Leftrightarrow3^{x+1}<1\Leftrightarrow x+1<0\Leftrightarrow x<-1.$
			Vậy tập nghiệm $S=(-\infty;-1)$.
		\end{itemchoice}
	}
\end{ex}

\begin{ex}%[Dự án chuẩn hóa 10-11, cấu trúc Bộ 2025]%[Lê Hồng Phi]%[1D7V2-4]
Trong hệ trục $Oxy$, cho hàm số $y=2x^2-3x$ và đường thẳng $\Delta\colon x+y-2=0$. Gọi $d$ là tiếp tuyến của đồ thị hàm số đã cho tại điểm có hoành độ bằng $x_0$. Các mệnh đề sau đúng hay sai?
\choiceTF
{Đường thẳng $\Delta$ có hệ số góc bằng $1$}
{\True Hệ số góc của $d$ là $k=4x_0-3$}
{Nếu $d\parallel \Delta$ thì $x_0=-1$}
{Nếu $d\parallel \Delta$ thì phương trình của $d$ là $d\colon 2x+2y-1=0$}
\loigiai{
\begin{itemchoice}
\itemch Phương trình đường thẳng $\Delta$ được viết lại là $y=-x+2$ và có hệ số góc bằng $-1$.
\itemch Ta có $y'=4x-3$ và hệ số góc của $d$ là $k=4x_0-3$.
\itemch Vì $d\parallel \Delta$ nên $4x_0-3=-1\Leftrightarrow 4x_0=2\Leftrightarrow x_0=\dfrac{1}{2}$.
\itemch Ta có $x_0=\dfrac{1}{2}$ nên tung độ của tiếp điểm là $y_0=2\cdot\left(\dfrac{1}{2}\right)^2-3\cdot\dfrac{1}{2}=-1$.\\ Phương trình tiếp tuyến $d$ là $y=-1\left(x-\dfrac{1}{2}\right)-1\Leftrightarrow y=-x-\dfrac{1}{2}\Leftrightarrow 2x+2y+1=0$.
\end{itemchoice}
}
\end{ex}

\begin{ex}%[Dự án chuẩn hóa 10-11, cấu trúc Bộ 2025]%[Lê Hồng Phi]%[1H8V1-3]
	Cho tứ diện $ABCD$ có $M$, $N$, $P$ lần lượt là trung điểm của $BC$, $AD$ và $AC$. Biết $AB=CD=a$, $MN=\dfrac{a\sqrt{3}}{2}$. Các mệnh đề sau đúng hay sai?
	\choiceTF
	{\True Số đo góc giữa hai đường thẳng $AB$ và $CD$ bằng góc giữa hai đường thẳng $PM$ và $PN$}
	{\True Tam giác $MPN$ là tam giác cân}
	{Công thức tính cô-sin góc $\widehat{MPN}$ là $\cos\widehat{MPN}=\dfrac{PM^2+PN^2-MN^2}{PM\cdot PN}$}
	{\True Số đo góc giữa hai đường thẳng $AB$ và $CD$ bằng $60^\circ$}
	\loigiai{

\begin{itemchoice}
\immini{
\itemch Ta có $PM$, $PN$ lần lượt là đường trung bình của các tam giác $ABC$ và $ACD$ nên $PM\parallel AB$, $PN\parallel CD$. Do đó $\left(AB,CD\right)=\left(PM,PN\right)$.
\itemch Ta có $PM$, $PN$ lần lượt là đường trung bình của các tam giác $ABC$ và $ACD$ nên $PM=\dfrac{1}{2}AB=\dfrac{1}{2}CD=PN$. Như thế, tam giác $MPN$ cân tại $P$.
\itemch Trong tam giác $MPN$ ta có $\cos\widehat{MPN}=\dfrac{PM^2+PN^2-MN^2}{2\cdot PM\cdot PN}$.
}{\begin{tikzpicture}[scale=1, font=\footnotesize, line join=round, line cap=round,>=stealth]
\path (0,0)coordinate (B) (4,0) coordinate (D) (1,-1) coordinate (C) ($(B)+(1.5,2.5)$) coordinate (A) ($(B)!0.5!(C)$) coordinate (M) ($(A)!0.5!(D)$) coordinate (N) ($(A)!0.5!(C)$) coordinate (P); 
\draw (A)--(B)--(C)--(D)--cycle (A)--(C) (M)--(P)--(N);
\draw[dashed] (B)--(D) (M)--(N);
\foreach \p/\g in {B/180,C/-90,D/0,A/90,M/225,N/30,P/135} \fill[black] (\p) circle(1pt)+(\g:0.3) node{$\p$};
\end{tikzpicture}}
\itemch Ta có $\cos\widehat{MPN}=\dfrac{PM^2+PN^2-MN^2}{2\cdot PM\cdot PN}=\dfrac{\left(\tfrac{a}{2}\right)^2+\left(\tfrac{a}{2}\right)^2-\left(\tfrac{a\sqrt{3}}{2}\right)^2}{2\cdot\tfrac{a}{2}\cdot\tfrac{a}{2}}=-\dfrac{1}{2}\Rightarrow\widehat{MPN}=120^\circ$.\\
Vậy $\left(AB,CD\right)=\left(PM,PN\right)=180^\circ-\widehat{MPN}=180^\circ-120^\circ=60^\circ$.
\end{itemchoice}
}
\end{ex}

\Closesolutionfile{ans}
\indapan{3}{ans/ans-cau-DS-1}

\caukq
\Opensolutionfile{ans}[ans/ans-cau-KQ-1]
\begin{ex}%[Dự án chuẩn hóa 10-11, cấu trúc Bộ 2025]%[Lê Hồng Phi]%[1D6H4-4]
Tìm nghiệm phương trình $\sqrt{2}\cdot 2^{3 x+1}=8$
\shortans{$0{,}5$} 
\loigiai{
Ta có $
\sqrt{2}\cdot2^{3x+1}=8 \Leftrightarrow 2^{3x+1}=4\sqrt{2} \Leftrightarrow 3x+1=\log\limits_24\sqrt{2}=\dfrac{5}{2} \Leftrightarrow x=\dfrac{1}{2}=0{,}5
$.\\
Vậy phương trình có nghiệm là $x=\dfrac{1}{2}$.
}
\end{ex}

\begin{ex}%[Dự án chuẩn hóa 10-11, cấu trúc Bộ 2025]%[Lê Hồng Phi]%[1D9V1-3]
Một bệnh truyền nhiễm có xác suất lây bệnh là $0{,}8$ nếu tiếp xúc với người bệnh mà không đeo khẩu trang; là $0{,}1$ nếu tiếp xúc với người bệnh mà có đeo khẩu trang. Chị Hoa có tiếp xúc với người bệnh hai lần, một lần đeo khẩu trang và một lần không đeo khẩu trang. Tính xác suất để chị Hoa bị lây bệnh từ người bệnh truyền nhiễm đó.
\shortans{$0{,}82$} 
\loigiai{
Gọi $A$ là biến cố: ``Chị Hoa bị nhiễm bệnh khi tiếp xúc người bệnh mà không đeo khẩu trang'' và $B$: ``Chị Hoa bị nhiễm bệnh khi tiếp xúc với người bệnh dù có đeo khẩu trang''.\\
 Dễ thấy $\overline{A}, \overline{B}$ là hai biến cố độc lập.\\
Xác suất để chị Hoa không nhiễm bệnh trong cả hai lần tiếp xúc với người bệnh là $\mathrm{P}(\overline{A}\cap\overline{B})=\mathrm{P}(\overline{A}) \cdot \mathrm{P}(\overline{B})=0{,}2 \cdot 0{,}9=0{,}18$.\\
Gọi $\mathrm{P}$ là xác suất để chị Hoa bị lây bệnh khi tiếp xúc người bệnh, ta có
\[\mathrm{P}=1-P(\overline{A}\cap\overline{B})=1-0{,}18=0{,}82.\]
}
\end{ex}

\begin{ex}%[Dự án chuẩn hóa 10-11, cấu trúc Bộ 2025]%[Lê Hồng Phi]%[1D6H1-1]
Biết $4^x+4^{-x}=23$, tính giá trị biểu thức $P=2^x+2^{-x}$.
\shortans{$5$} 
\loigiai{
Đặt $P=2^x+2^{-x}\Rightarrow P>0$.\\
Ta có $P^2=\left(2^x+2^{-x}\right)^2=4^x+4^{-x}+2 \cdot 2^x \cdot 2^{-x}=23+2=25$.\\
Vậy $P=5$.}
\end{ex}

\begin{ex}%[Dự án chuẩn hóa 10-11, cấu trúc Bộ 2025]%[Lê Hồng Phi]%[1D7H3-3]
	Một chất điểm chuyển động có phương trình $S=S(t)=t^3+4t^2-2$. Trong đó $t>0$, tính bằng giây $(\mathrm{s})$ và $S$ tính bằng $(\mathrm{m})$. Tính gia tốc của chuyển động tại thời điểm $t=2\mathrm{\,s}$.
\shortans{$20$ m/s$^2$} 
\loigiai{
Ta có vận tốc tức thời của chuyển động tại thời điểm $t$ là $v(t)=S'(t)=3t^2+8t$.\\
		Gia tốc tức thời của chuyển động tại thời điểm $t$ là $a(t)=v'(t)=6t+8$.\\
		Suy ra gia tốc của chuyển động tại thời điểm $t=2\mathrm{\,s}$ là $a(2)=6\cdot 2+8=20\mathrm{\,m}/\mathrm{s}^2$.
} 
\end{ex}

\begin{ex}%[Dự án chuẩn hóa 10-11, cấu trúc Bộ 2025]%[Lê Hồng Phi]%[1H8H4-3]
Cho lăng trụ đứng tam giác $ABC.A'B'C'$	có đáy $ABC$ là tam giác vuông cân tại $B$. Cho $AB=a$, $AA'=a\sqrt{3}$. Tính số đo góc giữa hai mặt phẳng $(A'BC)$ và $(ABC)$ theo đơn vị độ.
\shortans{$60^\circ$} 
\loigiai{
\immini{Vì lăng trụ đã cho là lăng trụ đứng và tam giác $(ABC)$ vuông ở $B$ nên $BC\perp A'B$.\\
Từ $(A'BC) \cap (ABC)=BC$, $AB\perp BC$, $A'B \perp BC$ suy ra góc giữa hai mặt phẳng $(A'BC)$ và $(ABC)$ bằng góc $\widehat{A'BA}=\varphi$.\\
Ta có $\tan\varphi=\dfrac{AA'}{AB}=\dfrac{a\sqrt{3}}{a}=\sqrt{3}$.\\
Vậy $\varphi=60^\circ$.
}{
\begin{tikzpicture}[scale=1, font=\footnotesize, line join=round, line cap=round,>=stealth]
\path (0,0)coordinate (A) (3,0) coordinate (C) (2,-1) coordinate (B) (0,3) coordinate (h);
\foreach \p in {A,B,C} {\draw (\p)--(\p)--+(h) coordinate (\p');}
\draw (A')--(B')--(C')--cycle (A)--(B)--(C) (A')--(B);
\draw[dashed] (A)--(C)--(A');
\foreach \p/\g in {A/180,B/-90,C/0,A'/180,B'/-30,C'/0} \fill[black] (\p) circle(1pt)+(\g:0.3) node{$\p$};
\foreach \x/\y/\z in {A/B/C,A'/A/B,A'/A/C}\draw pic[draw, angle radius=0.2cm]{right angle=\x--\y--\z};
\end{tikzpicture}
}
} 
\end{ex}

\begin{ex}%[Dự án chuẩn hóa 10-11, cấu trúc Bộ 2025]%[Lê Hồng Phi]%[1H8H5-3]
Cho hình chóp $S.ABCD$ có đáy là hình vuông cạnh bằng $2$, cạnh bên $SA$ vuông góc với đáy và $SA=2\sqrt{3}$. Tính khoảng cách từ điểm $B$ đến mặt phẳng $(SCD)$ (làm tròn kết quả đến $2$ chữ số sau dấu phẩy).
\shortans{$\approx 1{,}73$} 
\loigiai{
\immini{
Vì $AB$ song song với $(SCD)$ nên khoảng cách từ $B$ đến $(SCD)$ bằng khoảng cách từ $A$ đến $(SCD)$. Vẽ $AH$ vuông góc với $SD$ tại $H$. Ta có $CD\perp AD$ và $CD\perp SA$, suy ra $CD\perp AH$. Mặt khác $AH\perp SD$ nên suy ra $AH\perp (SCD)$. Do đó khoảng cách từ $A$ đến $(SCD)$ bằng độ dài $AH$. Ta có $$SD=\sqrt{SA^2+AD^2}=\sqrt{12+4}=4.$$ Từ hệ thức lượng $AH\cdot SD=SA\cdot AD$ ta tính được $$AH=\dfrac{SA\cdot AD}{SD}=\sqrt{3}.$$ Vậy khoảng cách cần tính bằng $\sqrt{3}\approx 1{,}73$.
}{\begin{tikzpicture}[scale=1, font=\footnotesize, line join=round, line cap=round,>=stealth]
\path (0,0)coordinate (B)++(3,0) coordinate (C)++(1.3,1) coordinate (D) ($(B)+(D)-(C)$) coordinate (A) (0,3) coordinate (h) ($(A)+(h)$) coordinate (S) ($(S)!0.6!(D)$) coordinate (H);
\draw (S)--(B)--(C)--(D)--cycle (S)--(C);
\draw[dashed] (S)--(A)--(B) (H)--(A)--(D);
\foreach \p/\g in {A/150,B/240,C/-60, D/0, S/90,H/45} \fill[black] (\p) circle(1pt)+(\g:0.3) node{$\p$};
\foreach \x/\y/\z in {S/A/B,S/A/D,S/H/A,A/B/C,B/C/D,C/D/A,D/A/B}\draw pic[draw, angle radius=0.2cm]{right angle=\x--\y--\z};
\end{tikzpicture}
}
} 
\end{ex}
\Closesolutionfile{ans}
\indapan{7}{ans/ans-cau-KQ-1}